\scanpage{336}
\mspage{37}

mais seulement au terme de la suivante. Pendant celle-ci, totalement libéré déjà des entraves de la vanité et de l'agressivité, vous aurez le rare privilège de servir alors pour la progression de tous les hommes plutôt que pour la vôtre, avec la coopération constante de Dieu et dans l'accomplissement d'une très grand mission qui vous aura été confiée par Lui.

14. La relation à nos parents domine notre vie

Dans ce que je viens de dire sur les racines de la vanité, il n'a été question que de l'aspect ``positif'' de la dépendance par rapport aux figures archétypes, à savoir du désir de leur plaire, ou celui de conquérir leurs faveurs. Ce désir conduit toujours à endosser des actes et des comportements qui ne naissent nullement des pulsions créatrices de l'être profond, mais de \textit{la soumission à des désirs supposés, toujours factices et toujours spirituellement stériles, que l'égo a été amené à prêter aux figures archétypes}, en réponse purement mécanique aux conditionnements reçus dans la présente existence et dans les existences précédentes. A dire vrai, il n'y a pas lieu de distinguer entre la Figure archétype elle-même, celle de la Mère disons, qui vit dans l'inconscient moyen, et cette \textit{image ``factice''} que l'égo se forme tout au cours du temps de l'errance, sous l'effet des divers conditionnements issus des circonstances de tout ordre qui entourent chacune de nos existences. Quant à l'archétype ``véritable'' de la Mère, on peut dire de lui qu'il est ``\textit{universel}'', c'est-à-dire qu'il est le même pour toutes les âmes. Il n'est la création ni de l'égo, ni de la psyché, ni même création de Dieu. Il existe en Dieu hors du temps, comme un des aspects éternels de Dieu : \textit{Dieu est Mère.} (Comme Il est aussi Père, comme Il est Epouse, comme Il est Epoux...) Et \textit{l'inconscient profond perçoit ``la Mère'' dans son essence véritable}, libre de toute déformation. C'est cette perception elle-même qui ``est'' la \textit{``véritable'' Figure archétype de la Mère, vivant, elle, dans l'inconscient profond et nourrissant ses pulsions créatrices tout au cours de nos existences} (tant terrestres, que dans l'au-delà). En elle réside la \textit{Réalité}, c'est-à-dire la perception véritable et fidèle que nous avons de la réalité. Cette Réalité est cachée plus que révélée par la \textit{Maya}, s'en donnant à coeur joie dans les couches moyennes et superficielles de la psyché, autour de l'image d'une ``Mère'' de pacotille (et parfois pas piquée de vers...).

Les conditionnements qui façonnent notre perception ``factice'' des Figures archétypes remontent presque tous à l'enfance. Ils sont déterminés en tout premier lieu par le milieu parental et par la relation des parents entre eux et à leur enfant (et, le cas échéant, à ses frères et soeurs). Sauf exceptions rarissimes, dues à des situations familiales exceptionnelles, \textit{les premiers êtres investis par les figures archétypes au cours d'une existence sont la mère, investie à la fois du rôle de} Mère \textit{et} d'Epouse, \textit{et le père, investi de même d'office du rôle de} Père \textit{et} d'Epoux. En l'absence de père ou de présence masculine en tenant lieu, la mère presque toujours se trouve investie également des rôles ardus (pour une mère...) et de Père, et d'Epoux. Et inversement dans le cas, de beaucoup plus rare, où le père doit suppléer à l'absence d'une mère défaillante.

Il faut se rappeler cependant que \textit{dès notre naissance}, en plus des \textit{Figures archétypes réelles} de la Mère, de l'Epouse, du Père, et de l'Epoux, qui vivent dans l'inconscient profond d'une existence éternelle et immuable, nous apportons déjà avec nous les \textit{figures archétypes factices} qui en sont un reflet déformé, et qui s'agitent dans l'inconscient moyen. Elles sont le produit surtout des conditionnements auxquels nous nous sommes trouvés soumis au cours de nos précédentes existences, et des réactions (le

% archive sheet number: 336
